% ===========================================================================
% ΕΝΔΕΙΚΤΙΚΕΣ ΛΥΣΕΙΣ ΚΡΙΤΗΡΙΩΝ ΑΞΙΟΛΟΓΗΣΗΣ — ΞΕΧΩΡΙΣΤΟ ΑΡΧΕΙΟ
% ===========================================================================
% Η αναφορά γίνεται με \label μέσα στο αντίστοιχο AssessmentQuestion.
% Δεν αντιγράφουμε ξανά την εκφώνηση. Αν αλλάξει η σειρά ερωτημάτων,
% το \ref ενημερώνεται στη δεύτερη μεταγλώττιση.
% Οι λύσεις εδώ είναι ενδεικτικές και όχι πλήρες λυσάρι του δοκιμίου.
\clearpage
\thispagestyle{plain}
\phantomsection
\addcontentsline{toc}{chapter}{Λύσεις κριτηρίων αξιολόγησης}
\BookRememberRunningChapter{Λύσεις κριτηρίων αξιολόγησης}%
\gdef\BookRunningChapterNumber{}%
\markboth{Λύσεις κριτηρίων}{Λύσεις κριτηρίων}
\BookSection{Λύσεις κριτηρίων αξιολόγησης}
\BookMinorHeading{Κριτήριο 1 — Ευθύγραμμη κίνηση}
\begin{AssessmentSolution}{asq:KRI1-A1}
Σωστή επιλογή: γ. Για σταθερή διανυσματική ταχύτητα, η επιτάχυνση
είναι μηδενική.
\end{AssessmentSolution}
\begin{AssessmentSolution}{asq:KRI1-G1}
Για ευθύγραμμη ομαλή κίνηση η εξίσωση θέσης είναι
\[
  x=x_0+\upsilon t=2+3t\quad (\mathrm{SI}).
\]
\end{AssessmentSolution}
